% -------------------------------------------------------------
\chapter{The General Structure}
% -------------------------------------------------------------
 
% -------------------------------------------------------------
\section{Two Levels of Testing}
% -------------------------------------------------------------
 
When Knuth describes empirical randomness testing\cite{knuth:1997-taocp-2}, two distinct levels are
at work, and it is important not to conflate them.
 
At the \textbf{upper level} sit the \emph{empirical tests} themselves ---
the frequency test, the serial test, the gap test, the poker test, the runs
test, and so on.  Each of these tests focuses on a specific property of the
draw sequence: how often each outcome appears, how pairs of consecutive
draws relate, how long the gaps between specific outcomes are, and so forth.
Applying one of these tests means extracting a particular numerical summary
from the sequence that reflects that property.
 
At the \textbf{lower level} sit two general-purpose statistical procedures:
the \emph{chi-squared test} and the \emph{Kolmogorov--Smirnov test}.  These
are not empirical tests in their own right.  They are the machinery that
every empirical test relies on to reach a conclusion.  Once an empirical
test has reduced the data to a test statistic, it is the chi-squared or KS
procedure that converts that statistic into a $p$-value and a verdict.
 
\begin{keyidea}[The Relationship Between the Two Levels]
Every empirical randomness test described by Knuth ultimately delegates its
final decision to one of two tools: the chi-squared test or the
Kolmogorov--Smirnov test.  These two procedures are not members of the
battery of empirical tests --- they are the statistical engine that all
members of the battery run on.
\end{keyidea}
 
% -------------------------------------------------------------
\section{The General Structure of a Statistical Test}
% -------------------------------------------------------------
 
Before examining the two procedures in detail, it is worth stating the
logical structure they share.
 
Each empirical test produces a \textbf{test statistic}: a single number
computed from the observed data.  This statistic is designed so that, when
the draw sequence truly is random, the statistic follows a \emph{known
reference distribution}.  The chi-squared and KS procedures each provide
one such reference distribution.
 
The conclusion of a test is then expressed as a \textbf{$p$-value}: the
probability that a genuinely random process would produce a test statistic
at least as extreme as the one observed.  A very small $p$-value means the
data is highly inconsistent with randomness; a $p$-value near one means the
data fits the expected pattern suspiciously well.  Both extremes are cause
for investigation.
 
We choose a \textbf{significance level} $\alpha$ in advance.  If the
$p$-value falls below $\alpha$, we reject the null hypothesis that the
sequence is random.  Common choices are $\alpha = 0.05$ and $\alpha = 0.01$.
 
% -------------------------------------------------------------
\section{What Makes These Two Procedures Universal}
% -------------------------------------------------------------
 
One might ask: why do all empirical tests converge on the same two
procedures?  The answer lies in a property that both the chi-squared and KS
tests possess, which makes them applicable across a wide variety of
situations without modification.
 
\begin{keyidea}[Distribution-Free Property]
The reference distributions of the chi-squared and Kolmogorov--Smirnov test
statistics under $H_0$ do not depend on the specific probability
distribution that the data is supposed to follow.  The same reference
distribution, and the same tables of critical values, apply regardless of
whether the underlying theoretical distribution is uniform, geometric,
exponential, or any other shape.
\end{keyidea}
 
This is what allows the same two procedures to serve as the foundation for
every empirical test.  Each empirical test examines a different property and
therefore works with a different theoretical distribution --- the frequency
test works with a uniform distribution over ball numbers, the gap test works
with a geometric distribution of waiting times, the runs test works with a
combinatorial distribution of run lengths, and so on.  Despite this variety,
every one of these tests can hand its test statistic to either the
chi-squared or the KS procedure and receive a valid $p$-value in return,
because those procedures are indifferent to the shape of the underlying
distribution.
 
\begin{important}
This distribution-free property is not trivial.  Most statistical tests are
derived specifically for one type of distribution and cannot be reused
elsewhere without re-derivation.  The chi-squared and KS tests are
exceptions: their validity is guaranteed by general mathematical theorems
that hold for any distribution, making them ideal as universal foundations.
\end{important}
 
% -------------------------------------------------------------
\section{The Chi-Squared Test}\label{sec:chi-squared-test}
% -------------------------------------------------------------
\index{Chi-squared Test}

\subsection{When It Is Used}
 
The chi-squared procedure is used when the test statistic from an empirical
test takes the form of \textbf{counts in discrete categories}.  The
empirical test observes how many times each of several possible outcomes
occurred and compares those counts to how many times each outcome was
expected to occur.
 
\subsection{The Test Statistic}
 
Let there be $k$ categories.  For category $i$, let $O_i$ denote the
observed count and $E_i$ the expected count under $H_0$.  The chi-squared
statistic is:
\[
  \chi^2 = \sum_{i=1}^{k} \frac{(O_i - E_i)^2}{E_i}
\]
 
The normalisation by $E_i$ is essential: it places discrepancies from
high-probability categories and low-probability categories on the same
scale, so that each contributes comparably to the total.
 
\subsection{The Reference Distribution and Its Independence from $E_i$}
 
When $H_0$ is true, the statistic $\chi^2$ follows approximately the
\emph{chi-squared distribution with $\nu = k - 1$ degrees of freedom}.
The critical observation is that this reference distribution depends only
on the number of categories $k$, and not at all on the individual expected
values $E_1, E_2, \ldots, E_k$.
 
This means that an empirical test with ten categories always consults the
chi-squared distribution with nine degrees of freedom, regardless of whether
those ten categories each have equal probability or wildly unequal
probabilities.  The same table entry provides the $p$-value in every case.
 
\begin{keyidea}
The chi-squared statistic translates category counts --- measured against
any theoretical distribution with any shape --- into a single universal
scale.  The reference distribution depends only on how many categories there
are, not on what those categories represent or how probable they are.
\end{keyidea}
 
\subsection{What the Chi-Squared Distribution Looks Like}

For each fixed number of degrees of freedom $\nu$, there is a corresponding
\emph{chi-squared distribution}.  You can think of it as a bell-shaped
curve that has been pushed over so that it starts at $0$ and has a long
tail out to the right:\index{Chi-squared Distribution}
\begin{itemize}
  \item $\chi^2$ can never be negative, because it is a sum of squared
        discrepancies.
  \item Values near $0$ correspond to counts that match the expectations
        almost perfectly.
  \item Larger values of $\chi^2$ correspond to counts that are
        increasingly uneven.
\end{itemize}

There is also a more structural way to view this distribution.  Take a
random variable $Z$ that follows a \emph{Gaussian} (or \emph{normal})
distribution with mean $0$ and variance $1$, usually written
$Z \sim \mathcal{N}(0,1)$.  This is the familiar bell-shaped curve: it is
symmetrical around $0$, values close to $0$ are common, and very large
positive or negative values are increasingly rare.  If you square such a
variable, $Z^2$, the result follows a chi-squared distribution with
one degree of freedom.  If you take several independent such variables
$Z_1, Z_2, \dots, Z_\nu$ and add their squares,
\[
  \chi^2 = Z_1^2 + Z_2^2 + \cdots + Z_\nu^2,
\]
then $\chi^2$ has a chi-squared distribution with $\nu$ degrees of
freedom.  In this sense, chi-squared distributions describe sums of
squared Gaussian variables.

For small $\nu$ the curve is strongly skewed; for large $\nu$ it becomes
more and more symmetric and bell-like.  The important point is that, once
$\nu$ is known, the entire shape of this curve is fixed, and so is the
probability that $\chi^2$ will fall in any given range.

\subsection{How the $p$-Value Is Obtained in Practice}

Once you have computed $\chi^2$ and the degrees of freedom $\nu$, the
$p$-value is simply the probability, under the chi-squared distribution
with $\nu$ degrees of freedom, of seeing a value at least as large as the
one you observed.  There are two practical ways to obtain it:
\begin{itemize}
  \item \textbf{Using tables.}  Traditional statistics books print
        \emph{chi-squared tables} with rows for degrees of freedom and
        columns for selected upper-tail probabilities (for example
        $0.10$, $0.05$, $0.01$, $0.001$).  Given $\nu$ and your $\chi^2$
        value, you look along the row for $\nu$ to see between which two
        columns your statistic falls.  This tells you that the $p$-value
        lies between the corresponding probabilities.
  \item \textbf{Using a computer program.}  Any scientific computing
        environment (R, Python, Julia, a spreadsheet, or a statistics
        package) can evaluate the chi-squared CDF directly.  You give it
        $\chi^2$ and $\nu$, and it returns the exact upper-tail
        probability.  This is the most convenient option in modern work,
        and is what we will assume when we quote numerical $p$-values in
        later chapters.
\end{itemize}

Either way, the logic is the same: we place the observed $\chi^2$ on the
chi-squared distribution with the correct degrees of freedom and read off
how much probability lies to the right of it.  That right-hand probability
is the $p$-value.
 
\subsection{What It Is Sensitive To}
 
The chi-squared test detects discrepancies that are \emph{distributed across
many categories}.  It is powerful when multiple categories deviate
moderately from their expected counts simultaneously.  Because the test
sums contributions from all categories equally, it does not weight any
particular category more heavily than another, and it discards any ordering
among the categories.  As a result, it cannot detect trends where deviations
are systematically arranged --- for example, consistently too high in one
part of the range and too low in another.
 
\subsection{A Requirement on Expected Counts}
 
The chi-squared distribution is an approximation, and it requires each
expected count $E_i$ to be sufficiently large to be reliable.  When an
empirical test produces categories with very small expected counts, those
categories must be merged with neighbouring ones until the requirement is
met.  This changes $k$ and therefore the degrees of freedom.
 
\subsection{How to Interpret the $p$-Value (Rule of Thumb)}
\index{$p$-value}

Knuth gives a practical interpretation scheme for the $p$-values produced by
randomness tests.\cite{knuth:1997-taocp-2}  The idea is not to treat a single
test run as a final verdict, but to classify outcomes by how extreme they are.

\begin{itemize}
  \item If $p < 1\%$ or $p > 99\%$, we \textbf{reject} the hypothesis that the
        numbers are sufficiently random.
  \item If $1\% \le p < 5\%$ or $95\% < p \le 99\%$, the numbers are
        \textbf{suspect}.
  \item If $5\% \le p < 10\%$ or $90\% < p \le 95\%$, the numbers are
        \textbf{almost suspect}.
\end{itemize}

Knuth also emphasises repetition: the chi-squared test is often performed at
least three times on different sets of data.  If at least two of the three
results are \emph{suspect}, the numbers are regarded as not sufficiently
random.\cite{knuth:1997-taocp-2}

% -------------------------------------------------------------
\section{The Kolmogorov--Smirnov Test}\label{sec:ks-test}
% -------------------------------------------------------------

\index{Kolmogorov--Smirnov Test}
\subsection{When It Is Used}
 
The KS procedure is used when the test statistic from an empirical test
takes the form of \textbf{ordered numerical values} that can be compared to
a continuous theoretical distribution.  Rather than grouping observations
into bins, it uses the full resolution of the data.
 
\subsection{The Empirical and Theoretical Cumulative Distribution Functions}
 
For any value $x$, the theoretical cumulative distribution function (CDF)
$F(x)$ gives the probability that a single observation falls at or below $x$
under $H_0$.  The empirical CDF $F_n(x)$, built from the $n$ observed values,
gives the fraction of observations that actually fell at or below $x$.
 
Sorting the observations as $U_{(1)} \leq U_{(2)} \leq \cdots \leq U_{(n)}$,
the KS test measures the largest vertical gap between these two curves.
Knuth defines the two one-sided statistics:
\[
  K^+ = \max_{1 \leq i \leq n} \left( \frac{i}{n} - F(U_{(i)}) \right)
\]
\[
  K^- = \max_{1 \leq i \leq n} \left( F(U_{(i)}) - \frac{i-1}{n} \right)
\]
 
$K^+$ is the largest point where the observed distribution exceeds the
theoretical one; $K^-$ is the largest point where it falls short.  Knuth
recommends reporting both separately, as they detect deviations in opposite
directions.
 
\subsection{The Reference Distribution and Its Independence from $F$}
 
When $H_0$ is true, the scaled statistics $K^+\sqrt{n}$ and $K^-\sqrt{n}$
follow the \emph{Kolmogorov distribution} --- a single tabulated reference
that does not depend on the form of $F$.
 
This universality follows from the \emph{probability integral transform}: if
$U$ is a random variable with continuous CDF $F$, then $F(U)$ is uniformly
distributed on $[0,1]$, for any $F$.  When the empirical test computes
$F(U_{(i)})$, the result is always a uniform sample regardless of the
original distribution $F$.  The KS test then compares this transformed
sample to the uniform distribution, which is always the same comparison,
making the reference distribution universal.
 
\begin{keyidea}
The KS procedure internally transforms the data so that it is always being
compared against a uniform distribution, no matter what the original
theoretical distribution $F$ is.  This transformation is exact and
assumption-free, making the reference distribution genuinely universal
rather than approximate.
\end{keyidea}
 
\subsection{What It Is Sensitive To}
 
The KS test detects discrepancies that are \emph{concentrated at a single
point} --- that is, a location where the empirical and theoretical CDFs
diverge most sharply.  Because it preserves the ordering of the data and
examines cumulative behaviour, it is sensitive to systematic trends: a
gradual drift where values are consistently above or below the theoretical
curve over a range.  Its focus on the single worst-case gap means it may
underweight many smaller but widespread deviations spread across the range.
 
Unlike the chi-squared test, the KS test requires no minimum expected count
and is exact for any sample size, making it particularly suitable when data
are sparse.
 
% -------------------------------------------------------------
\section{How the Two Procedures Complement Each Other}
% -------------------------------------------------------------
 
The chi-squared and KS tests are not redundant --- they are sensitive to
different shapes of departure from $H_0$, and for this reason the empirical
tests in Knuth's battery are designed to use whichever is better suited to
the particular test statistic at hand.
 
When an empirical test produces category counts, chi-squared is the natural
choice.  When it produces ordered numerical values that can be compared
directly to a theoretical CDF, KS is the natural choice.  Some empirical
tests can be formulated either way, allowing both procedures to be applied
and their conclusions compared.
 
\begin{keyidea}[Why Two Foundations Are Better Than One]
Chi-squared is powerful against distributed, category-wide discrepancies.
KS is powerful against concentrated, location-specific discrepancies.
Together, they cover the two most important shapes of departure from
randomness.  Because both are distribution-free, the choice between them
is driven purely by the form of the test statistic --- never by the shape
of the theoretical distribution being tested against.
\end{keyidea}
 
These two procedures are therefore not just tools of convenience.  They are
the mathematical foundation on which the entire structure of empirical
randomness testing rests.